%%
%% Automatically generated file from DocOnce source
%% (https://github.com/hplgit/doconce/)
%%
%%


%-------------------- begin preamble ----------------------

\documentclass[%
oneside,                 % oneside: electronic viewing, twoside: printing
final,                   % or draft (marks overfull hboxes, figures with paths)
12pt]{article}

\listfiles               % print all files needed to compile this document

\usepackage[a4paper, margin=1in]{geometry}

\usepackage{relsize,makeidx,color,setspace,amsmath,amsfonts}
\usepackage[table]{xcolor}
\usepackage{bm,microtype}

\usepackage[pdftex]{graphicx}

% Movies are handled by the href package
\newenvironment{doconce:movie}{}{}
\newcounter{doconce:movie:counter}


\usepackage{fancyvrb} % packages needed for verbatim environments
\usepackage{minted}
\usemintedstyle{default}

\usepackage[T1]{fontenc}
%\usepackage[latin1]{inputenc}
\usepackage{ucs}
\usepackage[utf8x]{inputenc}

% Set palatino as the default font family:
\usepackage[sc]{mathpazo}    % Palatino fonts
\linespread{1.05}            % Palatino needs extra line spread to look nice

%\usepackage{lmodern}         % Latin Modern fonts derived from Computer Modern

% Hyperlinks in PDF:
\definecolor{linkcolor}{rgb}{0,0,0.4}
\usepackage{hyperref}
\hypersetup{
    breaklinks=true,
    colorlinks=true,
    linkcolor=linkcolor,
    urlcolor=linkcolor,
    citecolor=black,
    filecolor=black,
    %filecolor=blue,
    pdfmenubar=true,
    pdftoolbar=true,
    bookmarksdepth=3   % Uncomment (and tweak) for PDF bookmarks with more levels than the TOC
    }
%\hyperbaseurl{}   % hyperlinks are relative to this root

\setcounter{tocdepth}{2}  % number chapter, section, subsection

% Tricks for having figures close to where they are defined:
% 1. define less restrictive rules for where to put figures
\setcounter{topnumber}{2}
\setcounter{bottomnumber}{2}
\setcounter{totalnumber}{4}
\renewcommand{\topfraction}{0.85}
\renewcommand{\bottomfraction}{0.85}
\renewcommand{\textfraction}{0.15}
\renewcommand{\floatpagefraction}{0.7}
% 2. ensure all figures are flushed before next section
\usepackage[section]{placeins}
% 3. enable begin{figure}[H] (often leads to ugly pagebreaks)
%\usepackage{float}\restylefloat{figure}

\usepackage[framemethod=TikZ]{mdframed}

% --- begin definitions of admonition environments ---

% Admonition style "mdfbox" is an oval colored box based on mdframed
% "notice" admon
\colorlet{mdfbox_notice_background}{yellow!30}
\newmdenv[
  skipabove=15pt,
  skipbelow=15pt,
  outerlinewidth=0,
  backgroundcolor=mdfbox_notice_background,
  linecolor=blue!80!black!20,
  linewidth=2pt,       % frame thickness
  frametitlebackgroundcolor=yellow!45,
  frametitlerule=true,
  frametitlefont=\normalfont\bfseries,
  shadow=false,        % frame shadow?
  shadowsize=11pt,
  leftmargin=0,
  rightmargin=0,
  roundcorner=5,
  needspace=0pt,
]{notice_mdfboxmdframed}

\newenvironment{notice_mdfboxadmon}[1][]{
\begin{notice_mdfboxmdframed}[frametitle=#1]
}
{
\end{notice_mdfboxmdframed}
}

% Admonition style "mdfbox" is an oval colored box based on mdframed
% "summary" admon
\colorlet{mdfbox_summary_background}{yellow!30}
\newmdenv[
  skipabove=15pt,
  skipbelow=15pt,
  outerlinewidth=0,
  backgroundcolor=mdfbox_summary_background,
  linecolor=blue!80!black!20,
  linewidth=2pt,       % frame thickness
  frametitlebackgroundcolor=yellow!45,
  frametitlerule=true,
  frametitlefont=\normalfont\bfseries,
  shadow=false,        % frame shadow?
  shadowsize=11pt,
  leftmargin=0,
  rightmargin=0,
  roundcorner=5,
  needspace=0pt,
]{summary_mdfboxmdframed}

\newenvironment{summary_mdfboxadmon}[1][]{
\begin{summary_mdfboxmdframed}[frametitle=#1]
}
{
\end{summary_mdfboxmdframed}
}

% Admonition style "mdfbox" is an oval colored box based on mdframed
% "warning" admon
\colorlet{mdfbox_warning_background}{yellow!30}
\newmdenv[
  skipabove=15pt,
  skipbelow=15pt,
  outerlinewidth=0,
  backgroundcolor=mdfbox_warning_background,
  linecolor=blue!80!black!20,
  linewidth=2pt,       % frame thickness
  frametitlebackgroundcolor=yellow!45,
  frametitlerule=true,
  frametitlefont=\normalfont\bfseries,
  shadow=false,        % frame shadow?
  shadowsize=11pt,
  leftmargin=0,
  rightmargin=0,
  roundcorner=5,
  needspace=0pt,
]{warning_mdfboxmdframed}

\newenvironment{warning_mdfboxadmon}[1][]{
\begin{warning_mdfboxmdframed}[frametitle=#1]
}
{
\end{warning_mdfboxmdframed}
}

% Admonition style "mdfbox" is an oval colored box based on mdframed
% "question" admon
\colorlet{mdfbox_question_background}{yellow!30}
\newmdenv[
  skipabove=15pt,
  skipbelow=15pt,
  outerlinewidth=0,
  backgroundcolor=mdfbox_question_background,
  linecolor=blue!80!black!20,
  linewidth=2pt,       % frame thickness
  frametitlebackgroundcolor=yellow!45,
  frametitlerule=true,
  frametitlefont=\normalfont\bfseries,
  shadow=false,        % frame shadow?
  shadowsize=11pt,
  leftmargin=0,
  rightmargin=0,
  roundcorner=5,
  needspace=0pt,
]{question_mdfboxmdframed}

\newenvironment{question_mdfboxadmon}[1][]{
\begin{question_mdfboxmdframed}[frametitle=#1]
}
{
\end{question_mdfboxmdframed}
}

% Admonition style "mdfbox" is an oval colored box based on mdframed
% "block" admon
\colorlet{mdfbox_block_background}{yellow!30}
\newmdenv[
  skipabove=15pt,
  skipbelow=15pt,
  outerlinewidth=0,
  backgroundcolor=mdfbox_block_background,
  linecolor=blue!80!black!20,
  linewidth=2pt,       % frame thickness
  frametitlebackgroundcolor=yellow!45,
  frametitlerule=true,
  frametitlefont=\normalfont\bfseries,
  shadow=false,        % frame shadow?
  shadowsize=11pt,
  leftmargin=0,
  rightmargin=0,
  roundcorner=5,
  needspace=0pt,
]{block_mdfboxmdframed}

\newenvironment{block_mdfboxadmon}[1][]{
\begin{block_mdfboxmdframed}[frametitle=#1]
}
{
\end{block_mdfboxmdframed}
}

% --- end of definitions of admonition environments ---

% prevent orhpans and widows
\clubpenalty = 10000
\widowpenalty = 10000

% --- end of standard preamble for documents ---


% insert custom LaTeX commands...

\raggedbottom
\makeindex

%-------------------- end preamble ----------------------

\begin{document}

% endif for #ifdef PREAMBLE


% ------------------- main content ----------------------



% ----------------- title -------------------------

\thispagestyle{empty}

\begin{center}
{\LARGE\bf
\begin{spacing}{1.25}
Liten JavaScript kurs rettet mot IT2
\end{spacing}
}
\end{center}

% ----------------- author(s) -------------------------

\begin{center}
{\bf Johan Hake${}^{}$ (\texttt{johan.hake@gmail.com})} \\ [0mm]
\end{center}

\begin{center}
% List of all institutions:
\end{center}
    
% ----------------- end author(s) -------------------------

\begin{center} % date
Aug 4, 2015
\end{center}

\vspace{1cm}


% inline figure
\centerline{\includegraphics[width=0.5\linewidth]{figures/javascript-seeklogo.pdf}}


\pagebreak

\section{Introduksjon}
\label{section:introduksjon}

Dette er en ressurs for elever og lærere som ønsker å bruke
javascript/HTML/CSS som programmeringsspråk i IT2-undervisningen. Hele
nettstedet/boka har betastatus, og innhold vil fylles på og endres
fortløpende. Dokumentet er basert på en WordPress
\href{{http://www.fuzzbin.org/IT2}}{\nolinkurl{version}} utarbeidet av Tom Jarle
Christiansen. Han har også lagd de fleste av kodeeksemplene og youtube
filmene. Et stort takk rettes derfor til Tom Jarle!

Relevante kodeeksempler finnes på GitHub og jsBin.

Nettstedet er bygget opp etter strukturen til kompetansemåla i faget,
(se avsnitt:~\ref{section:kompetansemål}), og inneholder teori, oppgaver
og videoer.

Bli gjerne med i nettverket for lærere som er interessert i IT2 med
javascript på \href{{https://www.facebook.com/groups/231993360276497}}{\nolinkurl{facebook}}.

Johan Hake, August 2015

% !split

\pagebreak

\section{Teori}

Den grunnleggende teorin i faget er delt opp i avsnitt som tilsvarer
kompetansemålene i faget:

\begin{itemize}
\item Multimedieutvikling, se avsnitt~\ref{section:teori:multimedieutvikling}

\item Programmering, se avsnitt~\ref{section:teori:programmering}

\item Planlegging og dokumentasjon, se avsnitt~\ref{section:teori:planlegging}
\end{itemize}

\noindent
% !split

\subsection{Multimedieutvikling}
\label{section:teori:multimedieutvikling}

% !split
\subsection{Programmering}
\label{section:teori:programmering}

% !split
\subsection{Planlegging og dokumentasjon}
\label{section:teori:planlegging}

% !split

\section{Læreplan i informasjonsteknologi}

Denne seksjonen er hentet fra \href{{http://www.udir.no/kl06/inf1-01/Hele}}{\nolinkurl{udir}}.

\subsection{Formål}

Informasjonsteknologien har hatt stor betydning for
samfunnsutviklingen de siste tiårene. Teknologien har i løpet av kort
tid endret kommunikasjonsmønsteret i samfunnet og skapt nye
arbeidsplasser og lærings- og forskningsarenaer. Samfunnet har behov
for mennesker som kan forstå, benytte og videreutvikle
informasjonsteknologien, men samfunnet trenger også mennesker med en
bevisst og kritisk holdning til hva teknologien gjør med mennesker og
samfunn. Informasjonsteknologien gir muligheter for å lage helt nye
produkter og tjenester gjennom kreativitet og samarbeid over
faggrenser, og bidrar dermed til teknologisk innovasjon.

Programfaget skal gi trening i kreativ tenkning og problemløsning og
i å formulere presise beskrivelser og finne generelle
mønstre. Programfaget skal bidra til å gi innsikt i hvordan
informasjon i form av tall, tekster, bilder, grafikk, film, lyd og
animasjoner kan struktureres og behandles automatisk som data, og
hvilke krav det setter til datamaskiner og annet digitalt
utstyr. Gjennom programfaget skal den enkelte få erfaring med bruk av
moderne teknologi og relevante utviklingsverktøy, og hvordan
sammensmelting av data-, lyd- og bildeteknologi kan gi rom for
skapende bruk av teknologien.

Opplæringen legger vekt på å konstruere IT-løsninger, og
informasjonsteknologi er derfor på mange måter et praktisk fag. Det
skal legges til rette for kommunikasjon og samarbeid i
programfaget. Programfaget informasjonsteknologi er et realfag, men
det har også sterke koblinger til mediefag, samfunnsfag, økonomi,
språkfag og formgivingsfag. Programfaget kan derfor gi et godt
grunnlag for studier innen ulike fagområder og for videreutvikling av
kompetanse i yrkeslivet.

\subsection{Grunnleggende ferdigheter}

Grunnleggende ferdigheter er integrert i kompetansemålene der hvor de
bidrar til utvikling av og er en del av fagkompetansen. I
informasjonsteknologi forstås grunnleggende ferdigheter slik:


\begin{block_mdfboxadmon}[]

\textbf{Å kunne uttrykke seg muntlig og skriftlig} i informasjonsteknologi
innebærer å planlegge og beskrive IT-løsninger, og utarbeide
brukerveiledninger og dokumentasjon. Videre vil det si å formulere
presise instruksjoner for datamaskiner i et programmeringsspråk. Det
betyr også å uttrykke seg på en klar og presis måte.

\textbf{Å kunne lese} i informasjonsteknologi innerærer å tolke beskrivelser,
brukerveiledninger, diagrammer, modeller, symboler og programkode på
en presis måte. Videre betyr det å forstå fagspesifikke tekster.

\textbf{Å kunne regne} i informasjonsteknologi innebærer å gjøre enkle
utregninger eller uttrykke formler i et programmeringsspråk. Videre
vil det si å bruke enkel matematisk logikk for å uttrykke en
betingelse på en presis måte.

\textbf{Å kunne bruke digitale verktøy} utgjør en grunnstamme i
informasjonsteknologi. Det innebærer å bruke IT-løsninger på en
effektiv måte og bruke digitale verktøy både i planleggings- og
dokumentasjonsprosesser.
\end{block_mdfboxadmon}



\subsection{Kompetansemål IT2}
\label{section:kompetansemål}

\paragraph{Planlegging og dokumentasjon.}
Hovedområdet handler om planlegging av IT-løsninger, og utvikling av
disse etter gitte spesifikasjoner for å oppfylle brukernes
behov. Videre dreier det seg om dokumentasjon og vurdering av
IT-løsninger. Hovedområdet omfatter også utforming, dokumentasjon og
vurdering av løsninger i forhold til retningslinjer for
brukergrensesnitt.


\begin{notice_mdfboxadmon}[Mål for opplæringen er at eleven skal kunne:]
\begin{enumerate}
\item spesifisere og begrunne funksjonelle krav til planlagte IT-løsninger 

\item velge og bruke relevante teknikker og verktøy for planlegging og utvikling av IT-løsninger

\item lage brukerveiledninger for IT-løsninger

\item gjøre rede for hvordan IT-løsninger utvikles i samarbeid mellom personer, og hvilke krav det setter til planleggings- og utviklingsprosessen

\item forklare hensikten med teknisk dokumentasjon og lage slik dokumentasjon for IT-løsninger, med spesiell vekt på å dokumentere grensesnitt mellom ulike delsystemer
\end{enumerate}

\noindent
\end{notice_mdfboxadmon}



\paragraph{Programmering.}
Hovedområdet handler om hvordan formelle språk kan brukes til å
formulere strukturer og sette sammen instruksjoner som kan utføres av
en datamaskin. Sentralt i hovedområdet er eksperimentering og
problemløsning. I tillegg dreier det seg om objektorientering.


\begin{notice_mdfboxadmon}[Mål for opplæringen er at eleven skal kunne:]
\begin{enumerate}
\item lese og bruke dokumentasjon og kode

\item definere variabler og velge hensiktsmessige datatyper

\item tilordne uttrykk til variabler

\item programmere med enkle og indekserte variabler eller andre kolleksjoner av variabler

\item programmere med valg og gjentakelser

\item lage egne og bruke egne og andres funksjoner eller metoder med parametere

\item programmere funksjoner eller metoder som blir aktivisert av hendelser

\item utvikle og sette sammen delprogrammer

\item teste og finne feil i programmer ved å bruke vanlige teknikker

\item gjøre rede for hensikten med objektorientert utvikling og begrepene klasse, objekt og arv
\end{enumerate}

\noindent
\end{notice_mdfboxadmon}

 

\paragraph{Multimedieutvikling.}
Hovedområdet handler om utforming, strukturering, implementering og
vurdering av multimedieapplikasjoner med tall, tekst, lyd, bilde,
video og animasjoner. I tillegg omfatter hovedområdet både
applikasjoner for lokal bruk og publisering over Internett.


\begin{notice_mdfboxadmon}[Mål for opplæringen er at elevene skal kunne:]
\begin{enumerate}
\item planlegge og utvikle multimedieapplikasjoner ved å kombinere egne og andres multimedieelementer av typene tekst, bilde, lyd, video og animasjoner

\item bruke programmeringsspråk i multimedieapplikasjoner

\item vurdere og bruke relevante filformater for tekst, bilde, lyd, video og animasjoner

\item vurdere multimedieprodukter med hensyn til brukergrensesnitt og funksjonalitet
\end{enumerate}

\noindent
\end{notice_mdfboxadmon}

 

% !split

% sphinx code-blocks: js=javascript pycod=python cod=fortran cppcod=c++ sys=console

\begin{minted}[fontsize=\fontsize{9pt}{9pt},linenos=false,mathescape,baselinestretch=1.0,fontfamily=tt,xleftmargin=7mm]{js}
var start_x = 80;
var start_y = 70;
var bredde = 60;
var hojde = 240;

rect(start_x, start_y, bredde, hojde);
rect(start_x+3*bredde, start_y, bredde, hojde);
rect(start_x+bredde, start_y+(hojde-bredde)/2, 2*bredde, bredde);

var c=document.getElementById("myCanvas");
var ctx=c.getContext("2d");

var start_x = 80;
var start_y = 70;
var bredde = 60;
var hojde = 240;

ctx.rect(start_x, start_y, bredde, hojde);
ctx.rect(start_x+3*bredde, start_y, bredde, hojde);
ctx.rect(start_x+bredde, start_y+(hojde-bredde)/2, 2*bredde, bredde);
ctx.stroke();



\end{minted}


\begin{notice_mdfboxadmon}[Kode eksempel:]
\begin{minted}[fontsize=\fontsize{9pt}{9pt},linenos=false,mathescape,baselinestretch=1.0,fontfamily=tt,xleftmargin=7mm]{html}
<!DOCTYPE html>
<html>
<head>
<meta name="description" content="Leke med div">
  <meta charset="utf-8">
  <title>JS Bin</title>
<style id="jsbin-css">
#minDiv{
 text-align: center;
	width: 200px;
	height: 100px;
	background-color: #aa5555;
	border-color: brown;
	border-width: 6px;
	border-style: solid;
	border-radius: 0px 40px;
	box-shadow: 10px 10px 5px #888888;
}
</style>
</head>
<body>
	<div id="minDiv"><h2>Jeg ligger i en div!</h2></div>


<script id="jsbin-source-css" type="text/css">#minDiv{
 text-align: center;
	width: 200px;
	height: 100px;
	background-color: #aa5555;
	border-color: brown;
	border-width: 6px;
	border-style: solid;
	border-radius: 0px 40px;
	box-shadow: 10px 10px 5px #888888;
}</script>
</body>
</html>
\end{minted}
\end{notice_mdfboxadmon}

 



\begin{doconce:movie}
\refstepcounter{doconce:movie:counter}
\begin{center}
\href{{http://www.youtube.com/embed/exjxJVQC3ms}}{\nolinkurl{http://www.youtube.com/embed/exjxJVQC3ms}}
\end{center}

\begin{center}  % movie caption
Movie \arabic{doconce:movie:counter}: For løkker
\end{center}
\end{doconce:movie}


% \href{{http://johanhake.github.io/js/js.do.txt}}{DocOnce source of this document}

% ------------------- end of main content ---------------


\printindex

\end{document}

